% Modelo reutilizavel para a disciplina Sistemas Embarcados com Arduino.
% Copie este arquivo para cada experimento e substitua os campos entre colchetes.
% Adapte secoes e atividades ao enunciado. Nao apresente exemplos como dados medidos.
\documentclass[12pt,a4paper]{article}

\usepackage[T1]{fontenc}
\usepackage[utf8]{inputenc}
\usepackage[brazil]{babel}
\usepackage{geometry}
\usepackage{graphicx}
\usepackage{float}
\usepackage{booktabs}
\usepackage{array}
\usepackage{listings}
\usepackage{xcolor}
\usepackage[hidelinks]{hyperref}

\geometry{left=3cm,right=2cm,top=3cm,bottom=2cm}
\setlength{\parindent}{1.25cm}
\setlength{\parskip}{0pt}

% Edite estes campos em cada novo relatorio.
\newcommand{\NumeroExperimento}{[numero]}
\newcommand{\TemaExperimento}{[tema do experimento]}
\newcommand{\AutorUm}{João Gabriel Lourenço}
\newcommand{\AutorDois}{[confirmar nome completo do segundo autor]}
\newcommand{\Cidade}{Marabá -- PA}
\newcommand{\AnoRelatorio}{2026}

\renewcommand{\lstlistingname}{Código}
\lstset{
  language=C++,
  basicstyle=\ttfamily\footnotesize,
  numbers=left,
  numberstyle=\tiny\color{gray},
  numbersep=8pt,
  frame=single,
  breaklines=true,
  breakatwhitespace=false,
  showstringspaces=false,
  tabsize=2,
  captionpos=b
}

\begin{document}

\begin{titlepage}
  \begin{center}
    Universidade Federal do Sul e Sudeste do Pará\\[0.5cm]
    Faculdade de Engenharia da Computação\\[0.5cm]
    Sistemas Embarcados com Arduino -- Prof. Dr. Elton Alves

    \vfill
    {\LARGE\bfseries Relatório do Experimento \NumeroExperimento\par}
    \vspace{0.7cm}
    {\Large \TemaExperimento\par}

    \vfill
    \AutorUm\\[0.4cm]
    \AutorDois

    \vfill
    \Cidade\\[0.4cm]
    \AnoRelatorio
  \end{center}
\end{titlepage}

\setcounter{page}{1}

\section{Introdução}
\textit{[Apresente o objetivo do experimento, o problema estudado e o que foi realmente realizado. Não copie resultados de relatórios anteriores.]}

\section{Materiais utilizados}
\begin{itemize}
  \item \textit{[Placa e componentes efetivamente usados.]}
  \item \textit{[Instrumentos, simulador e software, quando utilizados.]}
\end{itemize}

% Duplique a secao abaixo para cada projeto guiado, desafio ou tarefa avaliativa.
% Ajuste os titulos e subsecoes conforme o enunciado correspondente.
\section{[Projeto ou tarefa] -- [nome]}

\subsection{Objetivo e solução proposta}
\textit{[Descreva o requisito e a solução implementada. Distinga planejamento de funcionamento observado.]}

\subsection{Montagem}
\textit{[Descreva conexões, pinos, resistores, alimentação e demais componentes. Inclua uma imagem real ou esquema verificado quando disponível.]}

% Exemplo para inserir uma figura existente na mesma pasta do arquivo .tex:
% \begin{figure}[H]
%   \centering
%   \includegraphics[width=0.8\textwidth]{montagem.png}
%   \caption{Montagem da atividade.}
%   \label{fig:montagem-atividade}
% \end{figure}

\subsection{Código-fonte}
% Substitua pelo codigo efetivamente utilizado. Confira a versao no GitHub.
\begin{lstlisting}[caption={Código da atividade},label={lst:atividade}]
// Inserir aqui o codigo Arduino efetivamente utilizado.
\end{lstlisting}

\subsection{Explicação passo a passo}
\textit{[Explique as variáveis, setup(), loop(), entradas, saídas, decisões e temporização do código inserido.]}

\subsection{Histórico de implementação}
\textit{[Registre tentativas, erros, testes e soluções que ocorreram de fato. Se não houve dificuldade relevante, informe isso de modo objetivo.]}

% Se houver medicoes, substitua o exemplo abaixo por dados reais.
\subsection{Medições e análise}
\begin{table}[H]
  \centering
  \caption{Registros da atividade.}
  \label{tab:medicoes-atividade}
  \begin{tabular}{p{0.30\textwidth}p{0.30\textwidth}p{0.25\textwidth}}
    \toprule
    Condição & Leitura & Observação \\
    \midrule
    \textit{[preencher]} & \textit{[preencher]} & \textit{[preencher]} \\
    \bottomrule
  \end{tabular}
\end{table}
\textit{[Explique o método de obtenção dos dados e analise a relação entre as grandezas. Insira equações e gráfico quando o enunciado exigir.]}

\section{Resultados}
\textit{[Sintetize o que foi observado em todas as atividades e compare com os objetivos. Não declare como testado o que foi apenas simulado ou planejado.]}

\section{Conclusão}
\textit{[Retome os objetivos, os principais resultados verificados e as limitações. Não introduza dados novos.]}

\section{Vídeo das atividades}
\textit{[Insira o link verificável do vídeo exigido pelo enunciado, se aplicável.]}

\section{Declaração de uso de Inteligência Artificial}
\textit{[Informe o nível de uso de IA definido no enunciado da atividade e descreva com precisão como as ferramentas foram utilizadas. Não reutilize automaticamente a declaração de um relatório anterior.]}

\end{document}
